




\chapter{Introducción}
\label{ch:introduccion}
\addcontentsline{toc}{chapter}{Introducción}

\section{Contexto del secado solar de café en Panamá}
El café constituye uno de los productos agrícolas de mayor importancia económica y cultural para Panamá. Aunque el volumen de producción nacional es reducido en comparación con otros países productores, el café panameño ha alcanzado reconocimiento internacional gracias a la calidad de variedades como Geisha, obteniendo precios diferenciados en mercados especializados. En este contexto, preservar la calidad del grano durante las etapas posteriores a la cosecha resulta indispensable para mantener su competitividad y maximizar el valor agregado de la producción.\\

Una de las etapas más críticas del procesamiento del café es el secado, cuyo objetivo es reducir gradualmente el contenido de humedad del grano hasta niveles adecuados para su almacenamiento y comercialización. Un proceso de secado inadecuado puede provocar fermentaciones no deseadas, proliferación de hongos, aparición de defectos físicos y alteraciones en las propiedades sensoriales del café, afectando directamente su calidad y su valor comercial. Por esta razón, el control de la temperatura, la humedad y el tiempo de secado constituye un aspecto fundamental dentro de la cadena de producción.\\

Entre las diferentes tecnologías disponibles, los secadores solares tipo invernadero representan una alternativa de bajo costo operativo y reducido impacto ambiental, al aprovechar la radiación solar como fuente principal de energía térmica. Estos sistemas permiten disminuir el consumo de combustibles fósiles y ofrecen una solución particularmente atractiva para pequeños y medianos productores. Sin embargo, su desempeño depende en gran medida de las condiciones meteorológicas, tales como la irradiancia solar, la temperatura ambiente, la humedad relativa y la velocidad del viento, variables que cambian continuamente durante la operación y dificultan mantener condiciones internas de secado estables.\\

En los últimos años, el desarrollo de herramientas de simulación numérica, sistemas de adquisición de datos, sensores de bajo costo y técnicas de control inteligente ha abierto nuevas posibilidades para mejorar la eficiencia de los secadores solares. Estas tecnologías permiten comprender con mayor precisión los fenómenos de transferencia de calor y masa que ocurren en el interior del secador, facilitando el diseño de estrategias orientadas a incrementar la uniformidad del proceso, reducir los tiempos de secado y preservar la calidad final del café.\\


\section{Pertinencia institucional}
\subsection{Alineación con la Universidad Tecnológica de Panamá}
La Universidad Tecnológica de Panamá (UTP) orienta su misión hacia la formación integral de profesionales en ingeniería, ciencias y tecnología y la generación de conocimiento para impulsar el progreso del país, la transformación social y el desarrollo sostenible \parencite{utp2026}. En este contexto, la presente investigación se enmarca dentro de las líneas de trabajo de la Facultad de Ingeniería Eléctrica, al integrar conocimientos de automatización, control, instrumentación, inteligencia artificial y modelado computacional para resolver un problema de interés nacional.\\

El desarrollo de sistemas inteligentes para secadores solares de café constituye una aplicación multidisciplinaria de la ingeniería, en la que convergen áreas como la transferencia de calor y masa, la dinámica de fluidos computacional (CFD), el diseño de sistemas de adquisición de datos, la electrónica, el Internet de las Cosas (IoT) y los sistemas de control. Estas disciplinas forman parte de los campos de investigación y desarrollo promovidos por la Universidad Tecnológica de Panamá, orientados a generar soluciones tecnológicas con impacto económico, social y ambiental.\\

Asimismo, la investigación fortalece el vínculo entre la universidad y el sector productivo mediante la colaboración con el proyecto \" Coffee Solutions", el cual proporciona el caso de estudio utilizado para la caracterización experimental y la validación de los modelos desarrollados. Esta interacción favorece la transferencia de conocimiento desde la academia hacia la industria, promoviendo la aplicación de herramientas avanzadas de ingeniería para mejorar la competitividad del sector cafetalero nacional.\\

Finalmente, este trabajo contribuye a la formación de capacidades nacionales en áreas emergentes como la simulación numérica, la agricultura inteligente y los sistemas ciberfísicos, consolidando el papel de la Universidad Tecnológica de Panamá como una institución comprometida con la investigación aplicada y la innovación tecnológica al servicio del desarrollo del país.\\

\subsection{Alineación con el Plan Estratégico Nacional de Ciencia, Tecnología e Innovación (PENCYT 2025--2029)}

La presente investigación se encuentra alineada con el Plan Estratégico Nacional de Ciencia, Tecnología e Innovación (PENCYT 2025--2029), al proponer el desarrollo de tecnologías que contribuyen a la modernización del sector agroindustrial panameño mediante la integración de modelado computacional, inteligencia artificial, sistemas de control e instrumentación. En particular, el trabajo se relaciona con los siguientes vectores estratégicos del plan \parencite{senacyt2025}:

\begin{itemize}
\item \textbf{Vector Alimentación}, al fortalecer la productividad y sostenibilidad del sector cafetalero mediante tecnologías para la agricultura de precisión.

\item \textbf{Vector Digitalización}, mediante la aplicación de inteligencia artificial, Internet de las Cosas (IoT), simulación computacional y automatización.

\item \textbf{Vector Medioambiente}, al promover el aprovechamiento de la energía solar y el desarrollo de tecnologías que favorecen una producción más sostenible.

\item \textbf{Vector Transformación Productiva}, mediante la incorporación de sistemas inteligentes que incrementan la eficiencia y competitividad de los procesos agroindustriales.
\end{itemize}

\subsection{Contribución a los Objetivos de Desarrollo Sostenible}
La presente investigación se relaciona con los Objetivos de Desarrollo Sostenible (ODS) de la Agenda 2030 de las Naciones Unidas \parencite{onu2015}, al proponer una solución tecnológica orientada a mejorar la productividad agrícola mediante el uso de energías renovables, la automatización y la optimización de procesos. En particular, el trabajo se alinea con los siguientes objetivos:

\begin{itemize}
\item \textbf{ODS 2. Hambre Cero}, al contribuir al fortalecimiento de la producción agrícola mediante tecnologías para el procesamiento poscosecha.

\item \textbf{ODS 7. Energía Asequible y No Contaminante}, al promover el aprovechamiento de la energía solar como fuente renovable para el secado del café.

\item \textbf{ODS 9. Industria, Innovación e Infraestructura}, mediante el desarrollo de modelos computacionales, sistemas de control inteligente e instrumentación aplicada a la agroindustria.

\item \textbf{ODS 12. Producción y Consumo Responsables}, al favorecer procesos de secado más eficientes y con menor consumo de recursos.

\item \textbf{ODS 13. Acción por el Clima}, al fomentar tecnologías que reducen la dependencia de combustibles fósiles y disminuyen el impacto ambiental del procesamiento del café.
\end{itemize}


\section{Impacto esperado}
\subsection{Impacto socioeconómico}
La implementación de sistemas inteligentes para el control de secadores solares de café tiene el potencial de incrementar la competitividad del sector cafetalero panameño al mejorar la calidad y uniformidad del proceso de secado. Esto puede traducirse en una reducción de pérdidas poscosecha, una mayor estabilidad en la calidad del producto y un incremento en el valor comercial del café. Asimismo, el desarrollo de tecnologías de bajo costo favorece su adopción por pequeños y medianos productores, fortaleciendo la productividad y sostenibilidad del sector.

\subsection{Impacto ambiental}
La investigación promueve el aprovechamiento de la energía solar como fuente principal para el secado del café, reduciendo la dependencia de combustibles fósiles y las emisiones asociadas a los sistemas de secado convencionales. Adicionalmente, el uso de estrategias de control permite optimizar el consumo energético y mejorar la eficiencia del proceso, contribuyendo a una producción agrícola más sostenible.

\subsection{Impacto en igualdad e inclusión}
El diseño de una metodología basada en tecnologías de bajo costo, sensores comerciales y herramientas de software de código abierto busca facilitar el acceso a soluciones de automatización para productores con diferentes niveles de recursos. Esto contribuye a disminuir la brecha tecnológica en el sector agropecuario y promueve la transferencia de conocimiento hacia comunidades productoras, favoreciendo una mayor inclusión en la adopción de tecnologías agrícolas.


\section{Antecedentes del proyecto}
\subsection{El proyecto Coffee Solutions}
Coffee Solutions es una empresa panameña dedicada al desarrollo de soluciones tecnológicas para el secado solar de café. Como parte de sus iniciativas de innovación, la empresa desarrolló un secador solar tipo invernadero destinado a mejorar el proceso de secado del café mediante el aprovechamiento de la energía solar y un diseño constructivo orientado a preservar la calidad del grano.\\

Si bien el secador desarrollado constituye una alternativa sostenible frente a sistemas convencionales de secado, su operación continúa dependiendo en gran medida de las condiciones meteorológicas y del ajuste manual de sus variables de operación. Esta situación representa una oportunidad para incorporar herramientas de modelado, simulación y control automático que permitan mejorar el desempeño del sistema y reducir la variabilidad del proceso de secado.\\

En este contexto, Coffee Solutions proporciona la infraestructura experimental utilizada durante esta investigación, permitiendo la caracterización del secador, la adquisición de datos experimentales y la validación parcial de los modelos desarrollados.\\

\subsection{Alcance de la presente investigación}
La presente investigación no tiene como objetivo desarrollar una solución exclusiva para el secador solar de Coffee Solutions. Su propósito consiste en establecer una metodología integral para el modelado, simulación y diseño de sistemas inteligentes de control de temperatura y humedad aplicables a secadores solares de café tipo invernadero.\\

Para ello, se emplea el secador de Coffee Solutions como caso de estudio para la obtención de datos experimentales, la validación de modelos determinísticos y no determinísticos, y la evaluación de diferentes paradigmas de control mediante simulación. De esta manera, los resultados obtenidos buscan generar conocimiento y metodologías que puedan adaptarse a otros secadores solares con características similares, contribuyendo al desarrollo de tecnologías de automatización para el sector cafetalero.\\


\section{Justificación científica}
Aunque durante las últimas décadas se han desarrollado múltiples investigaciones sobre secadores solares de café, la mayoría se ha enfocado en aspectos específicos del problema, tales como el diseño geométrico, el análisis térmico, la dinámica de fluidos, la instrumentación o la implementación de estrategias particulares de control. En consecuencia, existe una limitada integración entre estas disciplinas dentro de una metodología sistemática que permita comprender el comportamiento completo del proceso de secado y facilite el diseño de sistemas de control inteligentes.\\

La presente investigación busca contribuir a cerrar esta brecha mediante el desarrollo de una metodología integral para el modelado y diseño de sistemas de control de temperatura y humedad en secadores solares de café tipo invernadero. Para ello, se integran modelos determinísticos basados en dinámica de fluidos computacional (CFD), modelos reducidos de transferencia de calor, técnicas de aprendizaje automático, validación experimental y el análisis de diferentes paradigmas de control dentro de un mismo marco metodológico.\\

Adicionalmente, la investigación incorpora el desarrollo de una plataforma experimental de adquisición de datos, así como procedimientos de validación estadística que permiten evaluar la capacidad predictiva de los modelos propuestos. Esta integración facilita el estudio del comportamiento térmico del secador, la comparación entre enfoques físicos y basados en datos, y la formulación de estrategias de control adaptables a distintas condiciones de operación.\\

En consecuencia, el principal aporte científico de este trabajo no radica únicamente en el estudio de un secador solar específico, sino en el establecimiento de una metodología reproducible para el modelado, simulación, validación y diseño de sistemas inteligentes de control aplicables a secadores solares de café con características similares.


\section{Definición del problema}
Los secadores solares de café tipo invernadero operan mediante el aprovechamiento de la radiación solar y de las condiciones ambientales para generar el calor necesario durante el proceso de secado. Sin embargo, la energía disponible y las condiciones internas del secador varían continuamente debido a cambios en la irradiancia solar, la temperatura ambiente, la humedad relativa, la velocidad del viento y la nubosidad. Como consecuencia, la evolución de la temperatura y la humedad en el interior del secador presenta un comportamiento variable y difícil de predecir.\\

Esta variabilidad limita la capacidad de mantener condiciones uniformes y estables durante el proceso. En particular, pueden formarse gradientes espaciales de temperatura y humedad que ocasionan que distintas zonas de la carga de café se sequen a velocidades diferentes. Estas diferencias pueden producir secado desigual, prolongación innecesaria del proceso, sobreexposición térmica de algunos granos y conservación excesiva de humedad en otros, con posibles efectos negativos sobre la calidad y estabilidad del producto.\\

El problema no consiste únicamente en medir valores promedio de temperatura y humedad, sino en representar adecuadamente su distribución espacial y temporal dentro del secador. Una cantidad reducida o una ubicación inadecuada de sensores puede no detectar zonas críticas, mientras que una red de instrumentación demasiado densa incrementa los costos de adquisición, instalación, mantenimiento y procesamiento de datos. Por tanto, es necesario determinar una configuración de sensado que proporcione información suficiente para caracterizar el sistema sin imponer una infraestructura económicamente desproporcionada.\\

A su vez, el diseño de un sistema de control requiere disponer de un modelo capaz de representar la respuesta del secador ante perturbaciones meteorológicas y acciones de control. Los modelos basados exclusivamente en principios físicos pueden requerir elevados recursos computacionales, una caracterización detallada de los materiales y condiciones de frontera, y tiempos de ejecución incompatibles con determinadas aplicaciones de control. Por otro lado, los modelos basados únicamente en datos dependen de la cantidad, calidad y representatividad de los registros experimentales disponibles, y pueden producir predicciones inconsistentes cuando operan fuera de las condiciones utilizadas durante su entrenamiento.\\

Existe, por consiguiente, la necesidad de evaluar de manera sistemática modelos determinísticos, modelos basados en datos y estrategias híbridas que combinen información física y experimental. Estos modelos deben ser suficientemente precisos para representar la dinámica térmica y de humedad del secador, pero también suficientemente eficientes para su posible utilización dentro de sistemas de control y supervisión.\\

Otra dificultad consiste en seleccionar las acciones de control y los actuadores necesarios para modificar las condiciones internas del secador. La incorporación indiscriminada de ventiladores, extractores, sistemas de calentamiento u otros dispositivos puede aumentar el consumo energético, la complejidad del sistema y la inversión inicial. En consecuencia, el problema de control debe abordarse considerando simultáneamente la estabilidad de las variables, la uniformidad espacial, el tiempo de respuesta, la eficiencia energética, el costo de implementación y la capacidad de adaptación frente a diferentes condiciones meteorológicas.\\

En este contexto, el problema central de la investigación se formula de la siguiente manera:
\begin{quote}
¿Cómo desarrollar y evaluar una metodología de modelado y diseño de control inteligente que permita representar, predecir y regular las distribuciones de temperatura y humedad en secadores solares de café tipo invernadero, considerando la variabilidad meteorológica, la disponibilidad limitada de datos experimentales y las restricciones de costo, instrumentación y capacidad computacional?
\end{quote}

Para estudiar este problema, el secador solar de Coffee Solutions se utiliza como caso de estudio y plataforma experimental. No obstante, la investigación no se limita al desarrollo de una solución particular para esta instalación, sino que busca establecer criterios y procedimientos que puedan ser adaptados a otros secadores solares de características semejantes.


\section{Hipótesis}
\label{intro:hipotesis}
La integración de modelos determinísticos, modelos basados en datos, instrumentación experimental y técnicas de validación permitirá desarrollar una metodología capaz de representar con suficiente precisión la dinámica de temperatura y humedad de un secador solar de café tipo invernadero, proporcionando una base técnica adecuada para el diseño y evaluación de estrategias de control inteligente.


\section{Preguntas de investigación}

La presente investigación busca responder las siguientes preguntas:

\begin{enumerate}

\item ¿Cuáles son las variables físicas, las perturbaciones ambientales y las características constructivas que determinan el comportamiento térmico e higrométrico de un secador solar de café tipo invernadero?

\item ¿Qué nivel de precisión pueden alcanzar los modelos determinísticos, los modelos basados en datos y los enfoques híbridos para representar la dinámica de temperatura y humedad del secador, considerando datos experimentales obtenidos en condiciones reales de operación?

\item ¿Qué estrategias de control inteligente pueden formularse a partir de los modelos desarrollados para regular las condiciones internas del secador frente a perturbaciones meteorológicas?

\item ¿En qué medida la metodología propuesta constituye una herramienta válida para el modelado, la simulación y la evaluación de sistemas de control aplicables a secadores solares de café con características similares?

\end{enumerate}

\section{Tipo de investigación}
La presente investigación corresponde principalmente a una \textbf{investigación aplicada}, ya que busca desarrollar una metodología para el modelado y diseño de sistemas inteligentes de control de temperatura y humedad en secadores solares de café, orientada a resolver un problema tecnológico concreto mediante la integración de herramientas de ingeniería y ciencias computacionales.\\

El trabajo combina fundamentos teóricos de transferencia de calor, mecánica de fluidos, psicrometría, teoría de control e inteligencia artificial con actividades de modelado matemático, simulación computacional, instrumentación experimental y validación mediante datos obtenidos en un secador solar operando en condiciones reales. El propósito no es únicamente comprender el comportamiento del sistema, sino generar procedimientos y herramientas que puedan emplearse para el análisis y diseño de sistemas de control aplicables a secadores solares de características similares.\\

Desde el punto de vista metodológico, la investigación adopta un \textbf{enfoque cuantitativo}, puesto que fundamenta sus resultados en modelos matemáticos, simulaciones numéricas, mediciones experimentales y análisis estadísticos para evaluar la capacidad predictiva de los modelos desarrollados y el desempeño de las estrategias de control propuestas.


\section{Objetivos}
\label{intro:objetivos}

\subsection{Objetivo general}
Desarrollar una metodología para el modelado y diseño de sistemas inteligentes de control de temperatura y humedad en secadores solares de café tipo invernadero, mediante la integración de modelos determinísticos, modelos basados en datos, instrumentación experimental y técnicas de validación, que permita evaluar diferentes estrategias de control bajo condiciones variables de operación.

\subsection{Objetivos específicos}
\begin{enumerate}

\item Caracterizar el comportamiento térmico e higrométrico de un secador solar de café tipo invernadero, identificando las principales variables de estado, perturbaciones ambientales y parámetros relevantes para su modelado.

\item Desarrollar modelos determinísticos y modelos basados en datos que representen la dinámica de temperatura y humedad del secador, incorporando los fenómenos de transferencia de calor, psicrometría y las mediciones experimentales disponibles.

\item Diseñar e implementar una plataforma experimental de adquisición de datos que permita obtener información para la calibración y validación de los modelos desarrollados.

\item Validar los modelos mediante la comparación con datos experimentales utilizando métricas de desempeño y herramientas de análisis estadístico.

\item Formular y evaluar diferentes paradigmas de control inteligente, utilizando los modelos desarrollados como base para analizar su desempeño frente a perturbaciones meteorológicas mediante simulación.

\item Analizar el alcance y las limitaciones de la metodología propuesta, identificando su potencial aplicación en secadores solares de café con características similares.

\end{enumerate}


%\section{Estructura del documento}
